In order to obtain the accurate estimation of the reconstruction efficiency detector readouts must be consistent and their responses in real data must be similar to those in MC. Therefore every dead/bad regions of every detector was removed from the analysis by introducing the well established procedure of fiducially cutting the dubious occupancy regions from the detector heatmaps. Fiducial cuts are presented on Fig. \ref{fig:dm_dce0}-\ref{fig:dm_emcalw3}.

Any additional mismatch between the real data and MC generated heatmaps (Fig. \ref{fig:acc_dce0} - \ref{fig:acc_emcalw3}) can be re-weighted and/or accounted as systematics (see Sec. \ref{sec:systematics_acceptance}). In this analysis inconsistencies were only found between real data and MC generated PC3w heatmaps due to bad/dead regions in MC that are not present in real data. To account for that additional fiducial cuts to the MC were added (Fig. \ref{fig:acc_pc3w_sim_fixed}). The re-weight for PC3w was calculated as the ratio of the amount of remaining real data after applying additional MC fiducial heatmaps to the amount of remaining real data after applying fiducial heatmaps in PC3w. The re-weight value was estimated to be 1.022.
